\section{Discussion}

A raw docking correlation map is shaped by a ligand-wide tendency to score favourably across many
targets. The same near-uniform leading direction appears
in all eight scoring outputs tested on fixed poses, and it carries most of the variance in both
large Vina matrices. After subtracting each ligand's mean, the target-relative geometry is much
more concentrated than any of the four transformation-matched nulls. The
null comparison matters because centring raises dimension by construction; without it, a rise in
participation ratio would not distinguish residual structure from an algebraic consequence of centring.

That residual geometry is not a universal protein map. It changes with the ligand support on which
it is estimated. Maps from the lower and upper molecular-weight quartiles agree weakly, whereas
size-matched splits that hold out whole chemical groups agree at about 0.99. Molecular weight is a
post hoc domain coordinate rather than an identified cause. The map should therefore be estimated on
a pilot sample drawn from the source-library domain rather than assumed to transfer unchanged.
The size-conditioned null also reproduces much of the residual edge ordering, especially on
Docking-44. Residual edges are therefore not evidence of biochemical cross-reactivity by themselves.

On the primary connectivity-de-overlapped PDSP graph, which includes only target pairs with at least ten shared
compounds having quantifiable $K_\text{i}$ at both targets, agreement rises from 0.044 for the raw
map to 0.295 for the residual map. The difference is unusual under unrestricted two-sided target
relabeling, but its within-family adjusted probability is unresolved. More importantly, residual
agreement falls to 0.076 when right-censored assay cells are restored and to 0.046 when the docking
map is rebuilt on chemically shifted complete-case support. The PDSP result is therefore a
support-sensitive observation, not confirmation of general performance beyond homology.

The other panels narrow the interpretation further. On the Novartis safety resource, residual
partial agreement after sequence, family and support control gives $p<0.05$ on one of three
endpoints under family-preserving permutation. On the frozen kinase endpoint, centred docking
approaches the sequence-identity baseline on ROC-AUC (0.710 versus 0.718) without passing it, and remains below it on average
precision (0.301 versus 0.426). KiRHub moves in the opposite direction. Across these panels, centring
does not reliably outperform sequence-based reasoning among close homologues.

The scorer comparison shows why the shared axis should be checked rather than assumed to be specific
to Vina. Atom-pair-count outputs carry a larger shared component than the interaction-fingerprint
output, with empirical term sums between them (Figure~\ref{fig:axis}). This ordering is consistent with contact-count
features loading more strongly on molecular bulk, but it is descriptive: the outputs span about
four related lineages and all were evaluated on retained Vina poses. It should guide an audit, not
be read as a causal ranking of scorer designs or a comparison of docking accuracy.

The audit is inexpensive. Two hundred randomly selected ligands from the source-library domain
recover the full residual map at Spearman 0.918 on DOCKSTRING-58 and 0.942 on Docking-44, and
chemical-group-held-out recovery gives similar values. In an exploratory compression analysis,
the five leading modes retain 0.517 of the residual trace; truncation raises the point estimates of
rank agreement and ROC-AUC on all four kinase panels. Average-precision point estimates rise on only two, and
the six-test sensitivity and cross-panel interpretation remain unresolved and post hoc. We did not
tune over neighbouring values of $k$, and do not claim that five modes are optimal. Participation
ratio counts directions; it does not tell us which directions will transfer to an experimental
endpoint.

Several within-ligand corrections produce similar point estimates on the safety panel. Heavy-atom
normalisation followed by centring reaches 0.449 on the reported-bound endpoint, the ordinal
transform 0.366 and standard centring 0.372. Unresolved pairwise contrasts do not establish that
these transforms are equivalent, and the kinase panels show that they are not interchangeable.
The exact preprocessing rule should therefore accompany every reported target map.

\subsection{What this study adds}

Ligand-size correction, per-target normalisation and principal-component summaries of docking
profiles are established ideas \cite{pan2003,vigers2004,carta2007,sepresa2009,casey2009,
luo2017,fukunishi2005,fukunishi2006,fukunishi2018}. The contribution here is the audit around those
steps. Holding poses fixed shows that the shared direction is not unique to one scoring program.
Transformation-matched nulls show that centring alone does not explain the residual concentration. External
pharmacology, sequence and family baselines test what the residual map does and does not add.
Calibration curves quantify the sampling cost, and the released command-line tool applies the matrix,
domain and recovery checks to a new ligand-by-target matrix.

A multi-target docking matrix should therefore be read in three stages: check for a near-uniform
component, compare the centred map with matched nulls, and test the residual geometry on the intended
ligand domain against a decision-relevant baseline. Without these checks, a target-correlation heat
map is a visual summary rather than evidence of cross-reactivity.

\subsection{Scope and limitations}

The eight-scorer comparison is descriptive. Its criteria were written before the non-Vina scores
were inspected, but that timing cannot be verified from the released package. All outputs use
retained Vina poses, so the comparison does not capture poses another function would select in its
own search. Several outputs also share features or training lineages, and no commercial scoring
function was tested.

The external analyses are post hoc. Endpoints, support thresholds, predictor fusion and
homology-preserving schemes were chosen after the data had been seen. PDSP paired-QAP probabilities
are two-sided, with Holm and Bonferroni adjustments over the 12 correlated rows. Those adjustments
do not cover the wider analytical search and do not repair the absence of prospective specification.
QAP respects the dependence of edges that share targets, but it does not carry forward the sampling
uncertainty of PDSP edges estimated from different compound sets.
PDSP values also pool heterogeneous studies and assay conditions, a limitation not removed by
median aggregation or rank correlation.

PDSP also has censoring and docking-support boundaries. Restoring right-censored cells at their
reported bounds weakens the association, while complete-case docking support largely removes the
raw-to-residual separation and reverses the fusion increment. Complete-case deletion changes the
chemical domain, so it is not a pure imputation comparison. The result therefore does not support
using residual docking as a general counterscreen selector. The remaining association concerns a
target-pair graph conditional on quantifiable $K_\text{i}$ values and the primary docking support.

Three negative results constrain the use case. On the kinase endpoint, centred docking approaches
but does not exceed sequence identity. On KiRHub, ROC-AUC moves from 0.751 to 0.744 and average
precision from 0.331 to 0.289. At the single-ligand level, top-5 target retrieval falls to 16.4\%
from 23.6\% after centring. The method describes relationships among target columns; it is not a
general way to rank targets for an individual compound.

Results under homology-preserving nulls vary by panel. On PDSP, neither the raw-to-residual
continuous difference nor the fusion increment has a Holm-adjusted two-sided probability below
0.05 under within-family relabelling. On the safety panel, residual partial agreement gives
$p<0.05$ on only one of three endpoints, and the
raw map also retains partial agreement. Within KLIFS groups, neither the continuous kinase
association nor its sequence-and-pocket-adjusted version gives $p<0.05$. Most evidence for an increment beyond receptor
relatedness is therefore weak and panel-dependent.

Absolute agreement with an experimental map is weak evidence on its own. A stable-hash nuisance
basis with mean out-of-fold $R^2$ near zero still reproduces geometry agreement between 0.10 and
0.35, a band that includes several observed values. Paired raw-versus-residual comparisons against
one fixed endpoint are more informative, but they remain conditional on the chosen panels.

The docking structures are inherited and not uniformly human: five of 43 receptor chains come from
rat, turkey, sheep or mouse, while the sequence baselines are human. Two affected targets, ADRB1 and
HTR3A, appear in the PDSP graph. Dropping them leaves a smaller 119-pair association, but this is a
sensitivity rather than a substitute for uniformly matched structures.

Finally, the two large docking matrices are separately generated but overlap. They share 582 exact
molecules, 804 connectivity blocks, 1{,}340 Bemis--Murcko scaffolds, nine mapped target identities and
three receptor structures. DAVIS, PKIS2 and KiRHub also represent different assay contexts on the
same 20 target labels rather than independent target-panel replications.

\section{Conclusions}

A shared ligand-wide score axis is the leading component in every scoring output tested and dominates
the two large raw Vina matrices. The centred map retains structured geometry that is inexpensive to
estimate, but that geometry belongs to a declared receptor panel, scoring pipeline and ligand library. External
pharmacology supports its use as an exploratory target-pair diagnostic on some diversified panels;
mixed safety and kinase results rule out a universal cross-reactivity map or compound-level ranking
rule. The practical recommendation is therefore not to replace the raw map but to audit both:
compute the row-centred map alongside it, test the residual against matched nulls, estimate it
on a pilot sample from the intended chemical domain, and let the comparison between the two
maps --- not the centred map by itself --- carry the interpretation. Which map is preferable is
panel-dependent: centred maps have higher target-pair agreement on the diversified
secondary-pharmacology graphs, but lower agreement on KiRHub and lower ligand-wise retrieval.
